%%%%%%%%%%%%%%%%
%%% gen 32 %%%
%%%%%%%%%%%%%%%%

\Chapter{32}

\Verse{1}
{
	\hE{וַיַּשְׁכֵּם לָבָן בַּבֹּקֶר וַיְנַשֵּׁק לְבָנָיו וְלִבְנוֹתָיו וַיְבָרֶךְ אֶתְהֶם וַיֵּלֶךְ וַיָּשׁב לָבָן לִמְקֹמוֹ׃}
}
{
	\eE{
		32 And Laban got up early in the morning and kissed his
		sons and his daughters and blessed them and went. And
		Laban went back to his place,
	}
}
{
}

\Verse{2}
{
	\hE{וְיַעֲקֹב הָלַךְ לְדַרְכּוֹ וַיִּפְגְּעוּ בוֹ מַלְאֲכֵי אֱלֹהִים׃}
}
{
	\eE{and Jacob went his way. And angels of {\God} came upon him.}
}
{
}

\Verse{3}
{
	\hE{וַיֹּאמֶר יַעֲקֹב כַּאֲשֶׁר רָאָם מַחֲנֵה אֱלֹהִים זֶה וַיִּקְרָא שֵׁם הַמָּקוֹם הַהוּא מַחֲנָיִם׃}
}
{
	\eE{
		And Jacob said when he saw them, “This is a camp of {\God},”
		and he called that place’s name Mahanaim.
	}
}
{
}

\Verse{4}
{
	\hJ{וַיִּשְׁלַח יַעֲקֹב מַלְאָכִים לְפָנָיו אֶל עֵשָׂו אָחִיו אַרְצָה שֵׂעִיר שְׂדֵה אֱדוֹם׃}
}
{
	\eJ{
		AND HE SENT And Jacob sent messengers ahead of him to
		Esau, his brother, to the land of Seir, the territory of
		Edom,
	}
}
{
}

\Verse{5}
{
	\hJ{וַיְצַו אֹתָם לֵאמֹר כֹּה תֹאמְרוּן לַאדֹנִי לְעֵשָׂו כֹּה אָמַר עַבְדְּךָ יַעֲקֹב עִם לָבָן גַּרְתִּי וָאֵחַר עַד עָתָּה׃}
}
{
	\eJ{
		and commanded them saying, “You shall say this: ‘To my
		lord, to Esau, your servant Jacob said this, “I’ve stayed
		with Laban and delayed until now,
	}
}
{
%	\footnote{
%		I’ve stayed with Laban. Rashi conveys a tradition that this means: “I
%		lived temporarily with him, but I learned nothing from his bad ways.”
%		But Jacob does in fact learn from Laban. The deceiver meets a deceiver.
%		He learns how it feels to receive an injustice, not just how it feels to
%		dish it out. He learns that it hurts. It is a step in his development as
%		a man. And it apparently brings about the honorable way in which he
%		behaves toward Esau when he returns.
%	}
}

\Verse{6}
{
	\hJ{וַיְהִי לִי שׁוֹר וַחֲמוֹר צֹאן וְעֶבֶד וְשִׁפְחָה וָאֶשְׁלְחָה לְהַגִּיד לַאדֹנִי לִמְצֹא חֵן בְּעֵינֶיךָ׃}
}
{
	\eJ{
		and ox and ass and sheep and male and female servant have
		become mine. And I’m sending to tell my lord so as to find
		favor in your eyes.”’”
	}
}
{
}

\Verse{7}
{
	\hJ{וַיָּשֻׁבוּ הַמַּלְאָכִים אֶל יַעֲקֹב לֵאמֹר בָּאנוּ אֶל אָחִיךָ אֶל עֵשָׂו וְגַם הֹלֵךְ לִקְרָאתְךָ וְאַרְבַּע מֵאוֹת אִישׁ עִמּוֹ׃}
}
{
	\eJ{
		And the messengers came back to Jacob saying, “We came to
		your brother, to Esau, and also he’s coming to you—and
		four hundred men with him.”
	}
}
{
}

\Verse{8}
{
	\hJ{וַיִּירָא יַעֲקֹב מְאֹד וַיֵּצֶר לוֹ וַיַּחַץ אֶת הָעָם אֲשֶׁר אִתּוֹ וְאֶת הַצֹּאן וְאֶת הַבָּקָר וְהַגְּמַלִּים לִשְׁנֵי מַחֲנוֹת׃}
}
{
	\eJ{
		And Jacob was very afraid, and he had anguish. And he
		divided the people who were with him and the sheep and the
		oxen and the camels into two camps.
	}
}
{
}

\Verse{9}
{
	\hJ{וַיֹּאמֶר אִם יָבוֹא עֵשָׂו אֶל הַמַּחֲנֶה הָאַחַת וְהִכָּהוּ וְהָיָה הַמַּחֲנֶה הַנִּשְׁאָר לִפְלֵיטָה׃}
}
{
	\eJ{
		And he said, “If Esau will come to one camp and strike it,
		then the camp that is left will survive.”
	}
}
{
}

\Verse{10}
{
	\hJ{וַיֹּאמֶר יַעֲקֹב אֱלֹהֵי אָבִי אַבְרָהָם וֵאלֹהֵי אָבִי יִצְחָק יְהֹוָה הָאֹמֵר אֵלַי שׁוּב לְאַרְצְךָ וּלְמוֹלַדְתְּךָ וְאֵיטִיבָה עִמָּךְ׃}
}
{
	\eJ{
		And Jacob said, “{\God} of my father Abraham and {\God} of my
		father Isaac, YHWH, who said to me, ‘Go back to your land
		and to your birthplace, and I’ll deal well with you,’
	}
}
{
}

\Verse{11}
{
	\hJ{קָטֹנְתִּי מִכֹּל הַחֲסָדִים וּמִכּל הָאֱמֶת אֲשֶׁר עָשִׂיתָ אֶת עַבְדֶּךָ כִּי בְמַקְלִי עָבַרְתִּי אֶת הַיַּרְדֵּן הַזֶּה וְעַתָּה הָיִיתִי לִשְׁנֵי מַחֲנוֹת׃}
}
{
	\eJ{
		I’m not worthy of all the kindnesses and all the
		faithfulness that you’ve done with your servant, because I
		crossed this Jordan with just my rod, and now I’ve become
		two camps.
	}
}
{
}

\Verse{12}
{
	\hJ{הַצִּילֵנִי נָא מִיַּד אָחִי מִיַּד עֵשָׂו כִּי יָרֵא אָנֹכִי אֹתוֹ פֶּן יָבוֹא וְהִכַּנִי אֵם עַל בָּנִים׃}
}
{
	\eJ{
		Save me from my brother’s hand, from Esau’s hand, because
		I fear him, in case he’ll come and strike me, mother with
		children.
	}
}
{
}

\Verse{13}
{
	\hJ{וְאַתָּה אָמַרְתָּ הֵיטֵב אֵיטִיב עִמָּךְ וְשַׂמְתִּי אֶת זַרְעֲךָ כְּחוֹל הַיָּם אֲשֶׁר לֹא יִסָּפֵר מֵרֹב׃}
}
{
	\eJ{
		And you’ ve said, ‘I’ll do well with you, and I’ll make
		your seed like the sand of the sea, that it won’t be
		countable because of its great number.’”
	}
}
{
}

\Verse{14}
{
	\hRJE{וַיָּלֶן שָׁם בַּלַּיְלָה הַהוּא וַיִּקַּח מִן הַבָּא בְיָדוֹ מִנְחָה לְעֵשָׂו אָחִיו׃}
}
{
	\eRJE{
		And he spent that night there. And he took an offering for
		Esau, his brother, from what had come into his hand:
	}
}
{
}

\Verse{15}
{
	\hE{עִזִּים מָאתַיִם וּתְיָשִׁים עֶשְׂרִים רְחֵלִים מָאתַיִם וְאֵילִים עֶשְׂרִים׃}
}
{
	\eE{
		two hundred she-goats and twenty he-goats, two hundred
		ewes and twenty rams,
	}
}
{
}

\Verse{16}
{
	\hE{גְּמַלִּים מֵינִיקוֹת וּבְנֵיהֶם שְׁלֹשִׁים פָּרוֹת אַרְבָּעִים וּפָרִים עֲשָׂרָה אֲתֹנֹת עֶשְׂרִים וַעְיָרִם עֲשָׂרָה׃}
}
{
	\eE{
		thirty nursing camels and their offspring, forty cows and
		ten bulls, twenty she-asses and ten he-asses.
	}
}
{
}

\Verse{17}
{
	\hE{וַיִּתֵּן בְּיַד עֲבָדָיו עֵדֶר עֵדֶר לְבַדּוֹ וַיֹּאמֶר אֶל עֲבָדָיו עִבְרוּ לְפָנַי וְרֶוַח תָּשִׂימוּ בֵּין עֵדֶר וּבֵין עֵדֶר׃}
}
{
	\eE{
		And he placed them in his servants’ hands, each herd by
		itself, and he said to his servants, “Pass on in front of
		me, and keep a distance between each herd and the next.”
	}
}
{
}

\Verse{18}
{
	\hE{וַיְצַו אֶת הָרִאשׁוֹן לֵאמֹר כִּי יִפְגשְׁךָ עֵשָׂו אָחִי וּשְׁאֵלְךָ לֵאמֹר לְמִי אַתָּה וְאָנָה תֵלֵךְ וּלְמִי אֵלֶּה לְפָנֶיךָ׃}
}
{
	\eE{
		And he commanded the first, saying, “When my brother,
		Esau, meets you and asks you, saying, ‘To whom do you
		belong, and where are you going, and to whom do these in
		front of you belong?’
	}
}
{
}

\Verse{19}
{
	\hE{וְאָמַרְתָּ לְעַבְדְּךָ לְיַעֲקֹב מִנְחָה הִוא שְׁלוּחָה לַאדֹנִי לְעֵשָׂו וְהִנֵּה גַם הוּא אַחֲרֵינוּ׃}
}
{
	\eE{
		then you’ll say, ‘To your servant, to Jacob. It’s an
		offering sent to my lord, to Esau. And here he is behind
		us as well.’”
	}
}
{
}

\Verse{20}
{
	\hE{וַיְצַו גַּם אֶת הַשֵּׁנִי גַּם אֶת הַשְּׁלִישִׁי גַּם אֶת כּל הַהֹלְכִים אַחֲרֵי הָעֲדָרִים לֵאמֹר כַּדָּבָר הַזֶּה תְּדַבְּרוּן אֶל עֵשָׂו בְּמֹצַאֲכֶם אֹתוֹ׃}
}
{
	\eE{
		And he also commanded the second, also the third, also all
		of those who were going behind the herds, saying, “You’ll
		speak this way to Esau when you find him,
	}
}
{
}

\Verse{21}
{
	\hE{וַאֲמַרְתֶּם גַּם הִנֵּה עַבְדְּךָ יַעֲקֹב אַחֲרֵינוּ כִּי אָמַר אֲכַפְּרָה פָנָיו בַּמִּנְחָה הַהֹלֶכֶת לְפָנָי וְאַחֲרֵי כֵן אֶרְאֶה פָנָיו אוּלַי יִשָּׂא פָנָי׃}
}
{
	\eE{
		and you’ll say, ‘Here is your servant, Jacob, behind us as
		well.’” Because he said, “Let me appease his face with the
		offering that’s going in front of me, and after that I’ll
		see his face; maybe he’ll raise my face.”
	}
}
{
%	\footnote{
%		his face … in front of me … I’ll see his face … my face … ahead of him.
%		The word “face” recurs five times in this line and all through these two
%		chapters (as does the word “raise”). It also is part of the term that is
%		translated as the preposition “in front of” or “ahead of” (in Hebrew,
%		lpny). The repetition conveys the force of this juncture in Jacob’s
%		life. He must face his past. He must face his brother, whom he wronged.
%		And in the middle of the account of his facing his brother will come the
%		account of his most immediate contact with God in his life, his struggle
%		after which he will say, “I’ve seen God face-to-face.” And the two
%		encounters, first with his God and then with his brother, will then be
%		brought together as he says to Esau, “I’ve seen your face—like seeing
%		God’s face!” (33:10). This is also ironic because Jacob left twenty
%		years earlier in the wake of his deception, which worked because his
%		weak-eyed father could not see his face!
%	}
}

\Verse{22}
{
	\hJ{וַתַּעֲבֹר הַמִּנְחָה עַל פָּנָיו וְהוּא לָן בַּלַּיְלָה הַהוּא בַּמַּחֲנֶה׃}
}
{
	\eJ{
		And the offering passed ahead of him. And he spent that
		night in the camp.
	}
}
{
}

\Verse{23}
{
	\hJ{וַיָּקם בַּלַּיְלָה הוּא וַיִּקַּח אֶת שְׁתֵּי נָשָׁיו וְאֶת שְׁתֵּי שִׁפְחֹתָיו וְאֶת אַחַד עָשָׂר יְלָדָיו וַיַּעֲבֹר אֵת מַעֲבַר יַבֹּק׃}
}
{
	\eJ{
		And he got up in that night and took his two wives and his
		two maids and his eleven boys and crossed the Jabbok ford.
	}
}
{
%	\footnote{
%		his eleven boys. His daughter, Dinah, is not mentioned.
%	}
}

\Verse{24}
{
	\hJ{וַיִּקָּחֵם וַיַּעֲבִרֵם אֶת הַנָּחַל וַיַּעֲבֵר אֶת אֲשֶׁר לוֹ׃}
}
{
	\eJ{
		And he took them and had them cross the wadi, and he had
		everything that was his cross.
	}
}
{
}

\Verse{25}
{
	\hJ{וַיִּוָּתֵר יַעֲקֹב לְבַדּוֹ וַיֵּאָבֵק אִישׁ עִמּוֹ עַד עֲלוֹת הַשָּׁחַר׃}
}
{
	\eJ{
		And Jacob was left by himself. And a man wrestled with him
		until the dawn’s rising.
	}
}
{
%	\footnote{
%		a man wrestled with him. With whom does he wrestle? It says “a man.” But
%		he is able against the man, and later the man names him yir-’l, which is
%		interpreted to mean “struggles with God,” and says to him, “You’ve
%		struggled with God and with people and were able.” And Jacob names the
%		place Peni-El, meaning “face of God,” and says, “I’ve seen God
%		face-to-face.” This all indicates that he has wrestled with God in human
%		form. But the prophet Hosea refers to this moment and says, “He fought
%		with God, and he fought with an angel and was able” (Hos 12:5–6). This
%		is consistent with the description of angels that I gave in the matter
%		of the three people who visit Abraham (18:3). Angels are material
%		expressions of God’s presence (emanations from the Godhead, hypostases).
%		When they are in front of a human, they look like a “man,” like
%		“people.” That is why Hosea can say “He fought with God” and then say in
%		poetic parallel about the same event “he fought with an angel.”
%	}
}

\Verse{26}
{
	\hJ{וַיַּרְא כִּי לֹא יָכֹל לוֹ וַיִּגַּע בְּכַף יְרֵכוֹ וַתֵּקַע כַּף יֶרֶךְ יַעֲקֹב בְּהֵאָבְקוֹ עִמּוֹ׃}
}
{
	\eJ{
		And he saw that he was not able against him, and he
		touched the inside of his thigh, and the inside of Jacob’s
		thigh was dislocated during his wrestling with him.
	}
}
{
%	\footnote{
%		dislocated. This is the term that is used in Num 25:4 for a punishment
%		that is inflicted on the Israelites following a heresy at Baal Peor. Its
%		meaning is uncertain both there and here. There it may mean to hang or
%		impale, in which case here it might mean that the thigh bone is hanging
%		out of joint.
%	}
}

\Verse{27}
{
	\hJ{וַיֹּאמֶר שַׁלְּחֵנִי כִּי עָלָה הַשָּׁחַר וַיֹּאמֶר לֹא אֲשַׁלֵּחֲךָ כִּי אִם בֵּרַכְתָּנִי׃}
}
{
	\eJ{
		And he said, “Let me go, because the dawn has risen.” And
		he said, “I won’t let you go unless you bless me.”
	}
}
{
%	\footnote{
%		I won’t let you go unless you bless me. Jacob demands a blessing. Is it
%		usual to demand a blessing from someone with whom we have been
%		fighting?! The struggle with God takes such an exertion of one’s
%		strength, and takes such a toll, that the human who does it has to live
%		with the wounds thereafter (Jacob limps after this struggle). And so he
%		is in need of blessing—and has earned it.
%	}
}

\Verse{28}
{
	\hJ{וַיֹּאמֶר אֵלָיו מַה שְּׁמֶךָ וַיֹּאמֶר יַעֲקֹב׃}
}
{
	\eJ{
		And he said to him, “What is your name?” And he said,
		“Jacob.”
	}
}
{
}

\Verse{29}
{
	\hJ{וַיֹּאמֶר לֹא יַעֲקֹב יֵאָמֵר עוֹד שִׁמְךָ כִּי אִם יִשְׂרָאֵל כִּי שָׂרִיתָ עִם אֱלֹהִים וְעִם אֲנָשִׁים וַתּוּכָל׃}
}
{
	\eJ{
		And he said, “Your name won’t be said ‘Jacob’ anymore but
		‘Israel,’ because you’ve struggled with {\God} and with
		people and were able.”
	}
}
{
%	\footnote{
%		Israel. There is little character development in Adam, Eve, Noah,
%		Abraham, Sarah, Isaac, or Rebekah, all of whom remain basically constant
%		figures through the stories about them. But Jacob changes, and the
%		matter of deception is intimately related to that development. As Esau
%		points out, Jacob’s very name connotes deception: to catch. And Jacob
%		starts out as a manipulator. But Jacob is changed after his experiences
%		in Mesopotamia. He has been the deceiver and the deceived. He has hurt
%		and been hurt. He is now a husband and a father, a man who has struggled
%		and prospered. For the rest of the story he is no longer pictured as a
%		man of action but, more often, as a relatively passive man, like his
%		father Isaac, seeking to appease his brother, avoiding strife and risk.
%		And precisely at the juncture that marks this change in Jacob’s
%		character he has his encounter with God at Penuel. God blesses him in a
%		remarkable etiology/etymology: “Your name shall no longer be called
%		Jacob, but Israel [yir-’l, understood here to mean ‘he struggles with
%		God’], because you’ve struggled with God and with people and were able.
%		Is this divine encounter the signpost of the change in Jacob’s
%		character, or the cause? Either way, as his character changes, and he
%		ceases to be the deceiver, just then he sheds the name Jacob (the one
%		who catches) and becomes instead Israel (the one who struggles with
%		God).
%	}
}

\Verse{30}
{
	\hJ{וַיִּשְׁאַל יַעֲקֹב וַיֹּאמֶר הַגִּידָה נָּא שְׁמֶךָ וַיֹּאמֶר לָמָּה זֶּה תִּשְׁאַל לִשְׁמִי וַיְבָרֶךְ אֹתוֹ שָׁם׃}
}
{
	\eJ{
		And Jacob asked, and he said, “Tell your name.” And he
		said, “Why is this that you ask my name?” And he blessed
		him there.
	}
}
{
}

\Verse{31}
{
	\hJ{וַיִּקְרָא יַעֲקֹב שֵׁם הַמָּקוֹם פְּנִיאֵל כִּי רָאִיתִי אֱלֹהִים פָּנִים אֶל פָּנִים וַתִּנָּצֵל נַפְשִׁי׃}
}
{
	\eJ{
		And Jacob called the place’s name Peni-El, “because I’ve
		seen {\God} face-to-face, and my life has been delivered.”
	}
}
{
}

\Verse{32}
{
	\hJ{וַיִּזְרַח לוֹ הַשֶּׁמֶשׁ כַּאֲשֶׁר עָבַר אֶת פְּנוּאֵל וְהוּא צֹלֵעַ עַל יְרֵכוֹ׃}
}
{
	\eJ{
		And the sun rose on him as he passed Penuel, and he was
		faltering on his thigh.
	}
}
{
}

\Verse{33}
{
	\hE{עַל כֵּן לֹא יֹאכְלוּ בְנֵי יִשְׂרָאֵל אֶת גִּיד הַנָּשֶׁה אֲשֶׁר עַל כַּף הַיָּרֵךְ עַד הַיּוֹם הַזֶּה כִּי נָגַע בְּכַף יֶרֶךְ יַעֲקֹב בְּגִיד הַנָּשֶׁה׃}
}
{
	\eE{
		On account of this the children of Israel to this day will
		not eat the tendon of the vein that is on the inside of
		the thigh, because he touched the inside of Jacob’s thigh,
		the tendon of the vein.
	}
}
{
}
